\section{Conclusion}
\label{sec:conclusion}

We presented a complete pipeline for adapting a general-purpose LLM to reliability engineering, a technical domain with no existing training datasets or benchmarks. Starting from 56 textbook-extracted seed questions, our adapted self-instruct methodology with cross-model answer verification and paraphrase augmentation scales the dataset to 866 verified examples. Supervised fine-tuning of Qwen3-8B yields a statistically significant $+18.6\%$ accuracy improvement ($p=0.002$) under 10-fold cross-validation, placing our 8B-parameter model between frontier models o3-mini and Gemini 3.1 Pro on domain-specific tasks.

Our analysis reveals two key findings for the synthetic data generation community. First, the \emph{self-instruct ceiling effect}: LLM generators create questions within their own capability envelope, making same-model answer consistency a trivially satisfied check rather than a meaningful quality filter. Second, paraphrase augmentation---creating surface-level rephrases while preserving mathematical content---is far more effective than generating harder synthetic questions, consistent with MetaMath~\cite{yu2024metamath} and suggesting that diversity of expression matters more than difficulty in small-data regimes.

GRPO with verifiable rewards achieves an additional $+7.5$pp on the 266 mixed-difficulty training questions, demonstrating that reinforcement learning can improve reasoning beyond SFT on verifiable tasks. Our ablation study reveals a paradoxical role for the SFT stage: while SFT alone causes catastrophic forgetting on unseen questions ($-20.4$pp on holdout), it is essential as initialization for GRPO ($+7.1$pp vs.\ GRPO from scratch). We hypothesize that SFT provides structural alignment (domain vocabulary, format, formula usage) that GRPO requires as a foundation for reasoning optimization.

Importantly, GRPO preserves base-model generalization: on the 54-question independent holdout, the GRPO model matches the base model (53.7\%), fully recovering the severe catastrophic forgetting caused by SFT alone ($-20.4$pp). This suggests that GRPO provides in-distribution gains without the out-of-distribution cost typically associated with fine-tuning on small datasets.

\paragraph{Limitations.} Our evaluation is limited to numeric questions with automated answer comparison; formula derivation and conceptual questions require human evaluation. The 54-question holdout is small (each question $\approx 1.85$pp), limiting statistical power for out-of-distribution conclusions. The self-instruct ceiling effect means our synthetic data may not cover the hardest problems in reliability engineering.

\paragraph{Future work.} Promising directions include: (1) scaling the GRPO training set via paraphrase augmentation of the 266 mixed questions, following the successful SFT scaling pattern; (2) introducing a KL penalty ($\beta > 0$) to reduce overfitting and preserve base model generalization; (3) early stopping based on holdout validation rather than training loss; (4) grounding data generation in textbook source material (following SearchInstruct~\cite{li2025searchinstruct}); and (5) extending to non-numeric question types and other engineering domains.
